\subsection{Applications}
\label{sec:applications}

This section zooms out to the downstream use-cases that actually justify an FTSG
model in a quantitative-finance workflow. We organize the discussion
along the task axis (FTSE, FTSI, FTSS) and end with a cross-cutting
paragraph on shared infrastructure (RL training gyms, risk-management
sandboxes, and synthetic-data sharing). Throughout, we restrict
ourselves to peer-reviewed or peer-reproducible works whose authors
explicitly evaluate against a real downstream application rather than
only on statistical fidelity.

\revnote{Secondary subheading ``Applications of FTSE'' removed; the paragraph below is now the lead-in of this topic within the primary section.}
\revadd{We begin with extrapolation (FTSE), the most mature sub-task in deployment.} The classical use-case
is to feed extrapolated point or quantile forecasts into a downstream
decision module, and the strongest evidence that this works comes from
the deep-reinforcement-learning (DRL) literature.
FinRL~\citep{liu2020finrl} and its successor
FinRL-Meta~\citep{liu2022finrlmeta} treat pretrained time-series
forecasters (LSTM, Transformer, or more recently LLM-based adapters) as
the observation pipeline for PPO/A2C/SAC trading agents trained on
Dow-30 and S\&P-500 constituents; their empirical study reports that
agents whose observation encoders are first fine-tuned on FTSE targets
achieve higher Sharpe and lower maximum drawdown than agents trained on
raw OHLCV. Hierarchical extensions such as
HRPM~\citep{wang2020hrpm} couple a portfolio-management FTSE head with
a separate order-execution head, explicitly accounting for price
slippage when large extrapolation errors are amplified by capital
allocation. More recently, LLM-driven forecasters
(Time-LLM~\citep{jin2024timellm}, CALF~\citep{liu2025calf}) have been
evaluated as the predictor component inside FinRL-style pipelines, with
consistent gains on A-share and crypto benchmarks.

A second, often-overlooked use of FTSE is volatility and risk
forecasting. TAIL-GAN~\citep{cont2025tailgan} is explicitly designed to
satisfy the joint elicitability of Value-at-Risk and Expected Shortfall,
and its authors plug the tail-respecting generator into a VaR back-test
to demonstrate regulatory-grade coverage. On the option-pricing side,
deep generative models of the underlying (e.g., quant GAN variants) have
been used to produce Monte-Carlo paths that replace historical
resampling for pricing exotic payoffs, where the FTSE tail behavior
dominates pricing error. FTSE outputs also feed
stochastic-portfolio-optimization engines: Hybrid CGAN for Portfolio
Analysis~\citep{lu2022hybrid} trains a conditional GAN whose label is
the Markowitz risk-aversion level, and the resulting synthetic future
paths are scored by an optimizer that selects weights without committing
to a single forecast.

\revnote{Secondary subheading ``Applications of FTSI'' removed.}
\revadd{Turning to imputation (FTSI), it is most often deployed in production at the data-engineering layer}, where its outputs feed classical pipelines rather than a learning agent. GAIN~\citep{yoon2018gain} introduced the generative-adversarial
imputation paradigm and remains a baseline in many open finance
datasets for handling missing quotes during off-hours, ticker
delistings, and exchange halts. Subsequent work has tightened the fit
between GAN imputation and financial time structure:
ImputeGAN~\citep{qin2023imputegan} and
AttnWGAIN~\citep{wang2025attn} replace generic CNN backbones with
Informer and diagonal-masked attention respectively, both reporting
smaller distributional drift on the K-line dataset than linear or MICE
baselines. Beyond price reconstruction, missing-data scenarios in
finance also arise at the macro layer.
FCI~\citep{herculano2025financial} uses Bayesian factor imputation to
keep Financial Condition Indices computable during incomplete-reporting
periods (e.g., during COVID-era reporting lags), demonstrating that
imputed macro factors can materially change the sign of the indicator
relative to linear interpolation.

A second family of applications concerns downstream classifier
performance on financial tabular data. For instance, evaluations of
Transformer-enhanced GAN oversampling~\citep{emaan2025improving} report
that synthesizing minority fraud transactions yields substantial gains
in detection recall compared to heuristic interpolation baselines such as
SMOTE on standard credit-card transaction benchmarks;
similarly, in credit-scoring ensembles~\citep{liu2021classification},
GAN-augmented minority classes materially reduce false-negative rates.
These studies reinforce that, in the financial setting, generative
imputation and augmentation are rarely the final product; rather, they
serve as value-adding pre-processing steps whose utility is judged by
downstream predictive metrics.

\revnote{Secondary subheading ``Applications of FTSS'' removed.}
\revadd{The broadest and most consequential deployment surface belongs to synthesis (FTSS). We organize it along the natural lifecycle of a quantitative strategy.}

\revadd{In strategy backtesting and replay,} ABIDES~\citep{byrd2019abides} and its Gym-compatible successor
ABIDES-Gym~\citep{amrouni2021abides} are agent-based discrete-event
simulators that wrap a synthetic limit-order-book generator in an
OpenAI Gym interface, enabling head-to-head backtests of RL agents that
interact with thousands of background traders.
\citep{balch2019how} compare single-agent market replay against
multi-agent interactive simulation and find that the interactive mode
systematically reorders the leaderboard of public trading algorithms
--- a striking demonstration that the choice of an FTSS simulation backbone can flip
strategy rankings. \citep{rollins2020which} extend this with a threaded
parallel simulator, scaling ABIDES-style order generation to high agent
counts, and reproduce qualitatively different ``pecking orders'' of
strategies compared with single-agent replay.

\revadd{For stress testing, scenario generation and risk management,} GAN- and diffusion-based synthesizers are increasingly used to
manufacture counterfactual or Black-Swan scenarios that historical
samples cannot supply. TAIL-GAN~\citep{cont2025tailgan} produces tail
scenarios that satisfy joint VaR/ES elicitability and are consumed by
bank-level stress-testing frameworks; the authors explicitly demonstrate
that the synthesized tails pass the Basel back-tests where naive
block-bootstrap fails. Hybrid CGAN for Portfolio
Analysis~\citep{lu2022hybrid} generates market scenarios conditioned on
risk aversion, allowing a portfolio engine to be re-evaluated against
thousands of hypothetical regimes in seconds. On the macro-finance
side, MacroSynth~\citep{alex2024macroeconomic} conditions on
macroeconomic factors to produce co-occurring crash and bull-run
scenarios, and the authors show that downstream mean-variance and
risk-parity allocations computed on the synthetic samples exhibit
sign-consistent risk premia relative to the historical record.
Diffusion-based generators are entering the same niche:
CoFinDiff~\citep{tanaka2025cofindiff} is evaluated as a scenario engine
for a portfolio whose weights are re-optimized on each generated
window, and Fin-Denoiser~\citep{wang2024denoiser} supplies high-SNR
scenarios for risk models that are notoriously sensitive to data
artifacts.

\revadd{As trading-agent training environments,} perhaps the most visible deployment of FTSS is as a training environment for RL
trading agents. FinRL-Meta~\citep{liu2022finrlmeta} ships near-real
market environments built from Yahoo Finance and CCXT feeds and
demonstrates that DRL agents pre-trained on these synthetic-augmented
environments generalize better to the held-out test year than agents
trained only on the historical window --- a result that is reproduced
across DOW-30, S\&P-500 and crypto-asset universes.
MarS~\citep{li2025mars} extends this idea to the order level by
generating asynchronous event streams, allowing agents to be evaluated
on metrics such as realized spread and adverse-selection cost that are
invisible to bar-level environments. Robust Risk-Sensitive RL
Agents~\citep{gao2021robust} stress-test the resulting agents by
perturbing the synthetic environment, showing that risk-aware training
reduces performance variance by an order of magnitude across random
seeds. Collectively, these works establish FTSS as the substrate for
``sim-to-real'' transfer in financial RL.

\revadd{In privacy-preserving data sharing and regulatory scenarios,} synthetic financial data is a recognized route to overcome the friction
between data sharing and the privacy constraints imposed by GDPR, CCPA,
and the equivalent Chinese PIPL. Six Levels of
Privacy~\citep{balch2024six} proposes a six-tier framework and
demonstrates empirically that synthetic equities, options and LOB
samples generated by TimeGAN, QuantGAN, and diffusion variants preserve
aggregate utility (Sharpe ratios, correlation matrices, tail-event
counts) while reducing membership-inference vulnerability to
near-baseline. Privacy Measurement in Tabular Synthetic
Data~\citep{boudewijn2023privacy} formalizes the corresponding
evaluation protocol that quantifies the privacy/utility frontier. On
the regulatory side, \citep{balch2024six} argue that FTSS pipelines can
produce synthetic benchmark datasets for third-party model validation
without exposing proprietary positions, a direction that has begun to
attract policy and supervisory attention in public-sector pilot studies.
This shift reframes FTSG not only as a research artifact but as
compliance infrastructure.

\revnote{Secondary subheading ``Cross-cutting Use: Synthetic Datasets as Benchmarks'' removed.}
\revadd{Beyond the per-task applications above, FTSG has produced a small but rapidly maturing ecosystem of benchmark datasets that are themselves reusable scientific infrastructure.} TSGBench~\citep{ang2023tsgbench}
provides paired real/synthetic time-series collections for GOOGL and FX
series and is consumed by downstream papers as a standardized fidelity
benchmark. CTBench~\citep{ang2025ctbench} extends this to Binance
cryptocurrency markets, where extreme volatility makes synthetic-data
quality both harder to achieve and more consequential.
FinTSB~\citep{hu2025fintsb} and FinTSBridge~\citep{wang2025fintsbridge}
package millions of Chinese A-share, NYSE and NASDAQ trajectories
together with pre-defined train/valid/test splits, lowering the cost of
entry for new FTSG research and making FTSE/FTSI results directly
comparable. These benchmarks invert the typical deployment pipeline:
rather than synthesizing data to fit a model, the model is now fitted
to a benchmark whose provenance, statistics and license are publicly
auditable. This is a subtle but important shift for reproducibility and
regulatory acceptance.

In summary, FTSE, FTSI, and FTSS are not merely three academic
formulations of the same underlying problem. Each maps onto a
distinct, value-generating workflow in industry. FTSE plugs into
decision engines (trading agents, portfolio optimizers, risk models);
FTSI into data-engineering layers (data cleaning, augmentation for
tabular classifiers); and FTSS into simulation infrastructure (RL
training, stress testing, privacy-preserving sharing). The applications
in this section collectively motivate the breadth of design choices
catalogued in the taxonomy and task-specific generation \revadd{sections}.

